\section{How the proof goes}\label{sec:arch}

The argument is Light's chain. This section describes each link and what had to be done to make
it hold at $\beta=0.96$. Formal names are given at each step; the emphasis is on the economics.

\subsection{Theorem 1: saving rises with the rate}\label{sec:thm1}

Fix two admissible rates $r_1<r_2$ and write $g_i$ for the policy at $r_i$. Light's Theorem 1 is
$g_1(a,z)\le g_2(a,z)$ for every state. His proof compares the two Bellman operators along
value-function iterates: the map $v\mapsto v(ta,\cdot)$, $t=R_1/R_2$, turns a household at the
higher rate with assets $ta$ into one with the same cash on hand as a household at the lower rate
with assets $a$, and the induction shows that the difference of the two value functions has
increasing differences in assets, which is what makes the higher-rate policy larger. The step where
the utility function enters is an inequality of the form $t\,u'(c(ta))\le u'(c(a))$ for the
rescaled consumption, and this is where relative risk aversion at most one is used: it is
monotonicity of $c\,u'(c)$, which for log utility is an identity.
(\lean{monotoneOn\_bellman\_sub\_scaled\_of\_piv}, \lean{policy\_mono\_withRate\_crra}.)

The induction needs the constraint slack at the compared states and the cap slack everywhere along
the iterates, and the published tests for both scale like $1/(1-\beta R)$. At Aiyagari's rates
$1-\beta R$ is about $0.002$, and the tests fail by orders of magnitude. What replaces them is in
the next two subsections. With those in hand the induction is run locally, on an interval of rates
short enough that the constants survive, and the local statements are chained across
$[r_{lo},r_{hi}]$ (\lean{LargeBetaLight}).

\subsection{The corner from the Euler inequality}\label{sec:corner}

\begin{lemma}[Nobody at the constraint consumes less than the lowest income]
At an interior choice $g(a,z)>0$ the Euler inequality reads $u'(c)\le\beta R\,\E[u'(c')]$. Take the
income state at which consumption at zero assets is smallest. If that household saved, tomorrow's
consumption would be at least today's minimum, so $u'(c)\le\beta R\,u'(c)$, impossible with
$\beta R<1$. So it is at the corner, consumes its income, and consumption at the constraint is at
least $y_{\min}$ in every state.
\formal{IncomeFluctuation.minIncome\_le\_consumptionFn\_floor}
\end{lemma}

The corner condition at a state $(a,z)$ is then the contrapositive of the Euler inequality: if
$\beta R\,u'(y_{\min})<u'(m(a,z))$, saving is not optimal, because saving would give
$u'(m(a,z))\le u'(c)\le\beta R\,\E[u'(c')]\le\beta R\,u'(y_{\min})$. This is a primitive
inequality, checkable by hand at $\beta=0.96$, and it is what
\lean{aiyagari1994\_log\_equilibriumRate\_unique} evaluates internally with an explicit corner
level $a_0=\min\{\bar a,\;y_{\min}(1-\beta R)/(2\beta R^2)\}$.
(\lean{EulerCorner}.)

\subsection{The cap from the minimal marginal propensity to consume}\label{sec:mpc}

\begin{lemma}[Minimal MPC]
For CRRA utility with parameter $\gamma$ at a zero borrowing limit, with $\beta R<1$,
\[
  c(a,z)\ \ge\ \kappa\, m(a,z),\qquad \kappa=1-\frac{\Thorn}{R},\quad \Thorn=(\beta R)^{1/\gamma}.
\]
\formal{IncomeFluctuation.crra\_minMPC\_mul\_le\_consumptionFn}
\end{lemma}

$\Thorn$ is Carroll's absolute patience factor and $\kappa$ the MPC of a perfect-foresight consumer
with no human wealth \citep{CarrollShanker2026}. The bound is self-reproducing under the Bellman
operator: if the continuation's consumption satisfies it, then at an interior choice the Euler
inequality gives $c^{-\gamma}\le\beta R(\kappa R a')^{-\gamma}$, hence $c\ge(\kappa R/\Thorn)a'$, and
$(\kappa R/\Thorn)/(1+\kappa R/\Thorn)=\kappa$ exactly because $1-\kappa=\Thorn/R$; at the corner
the household consumes everything. For log utility $\kappa=1-\beta$ at every rate, which is why the
log results carry no rate-dependent constant.

Two things follow. Saving is at most $(1-\kappa)m(a,z)$, so at the cap
$g(\bar a,z)\le\beta(y_{\max}+R\bar a)<\bar a$ whenever \eqref{eq:cap} holds, and the cap is slack
along every iterate and at the fixed point (\lean{crra\_policy\_lt\_cap\_of\_minMPC}). And
integrating against a stationary distribution, $K\le\beta(\bar Y+RK)$, so
\begin{equation}\label{eq:ceiling}
  K(\mu)\le\frac{\beta\,\bar Y}{1-\beta R},
\end{equation}
the log capital ceiling used in Theorems~\ref{thm:below} and \ref{thm:exist}
(\lean{log\_aggregateCapital\_le}, \lean{log\_aggregateCapital\_le\_mean}).

\subsection{Theorem 2: the coupling}\label{sec:coupling}

\begin{proposition}[Aggregate capital rises with the rate]
Let $g_1\le g_2$ pointwise on the region, and let the two Markov operators be run from a common
initial distribution $\mu_0$, converging to stationary distributions $\mu_1$ and $\mu_2$. Then
$\mu_2$ dominates $\mu_1$ in the asset coordinate, and $K(\mu_1)\le K(\mu_2)$.
\formal{IncomeFluctuation.monotoneOn\_capitalSupply\_of\_policy\_mono}
\end{proposition}

Dominance is tested against bounded continuous functions that are increasing in assets at each
income state. Three facts carry the argument: the Markov operator preserves that test class, because
the policy is increasing in assets (\lean{policy\_mono}); a pointwise-larger policy gives a
dominating operator on that class; and dominance is closed under weak limits. The asset coordinate
itself is in the test class, so the mean is ordered. This is the argument of \citet{Light2022},
Theorem 2, which asks for monotonicity in one coordinate only; monotonicity in the whole state
fails here because the policy need not rise with the income shock.

\subsection{Uniqueness and convergence of the stationary distribution}\label{sec:doeblin}

The coupling needs the iterates from $\mu_0$ to converge, and the equilibrium condition needs
capital supply to be a function of the rate, not a correspondence. Both come from a Doeblin
condition. The asset transition is deterministic given the shock, so two households with different
assets never meet, and no minorisation against a spread-out measure holds. The borrowing constraint
supplies an atom: if a run of $N$ draws of the worst income state $z_0$ drives every household to
exactly zero assets, then every $N$-step transition charges the single point $(0,z_0)$ with
probability at least $\Pi_{\min}(\cdot,z_0)^N$, which is a Doeblin condition. The contraction is run
on the function side: the oscillation of $T^N h$ is at most $(1-\epsilon)$ times that of $h$, so
two stationary distributions agree on every bounded continuous function, and the iterates from any
start converge weakly to the stationary distribution.
(\lean{stationary\_unique}, \lean{tendsto\_pushProb\_iterate}.)

That a run of bad shocks reaches the constraint is the \emph{decline} property: above some asset
level the household dissaves in the worst state. For an impatient household this is
\citet{Acikgoz2018}, Proposition 4, and it is proved here from the corner condition and the minimal
MPC (\lean{ImpatientDecline}). \ref{E2} is what makes the run have positive probability from every
state. At each rate in $[r_{lo},r_{hi}]$ the conclusion is a unique stationary distribution,
reached from every start (\lean{log\_existsUnique\_isStationary}).

\subsection{Theorem 3: single crossing, and Aiyagari's form of the condition}\label{sec:homog}

With supply a monotone function of the rate on $[r_{lo},r_{hi}]$ and $D$ strictly decreasing, two
equilibrium rates coincide: this is \lean{eq\_of\_monotoneOn\_of\_strictAntiOn} and needs no
continuity. The remaining step is to connect \eqref{eq:eqm} to Aiyagari's condition at the
equilibrium wage.

\begin{proposition}[Homogeneity of the stationary distribution]
Scaling every income and the cap by $w>0$ scales the policy by $w$. The map $(a,z)\mapsto(wa,z)$
therefore commutes with the Markov operator, so it carries stationary distributions of the unit-wage
economy to stationary distributions of the wage-$w$ economy, and aggregate capital scales by $w$.
\formal{IncomeFluctuation.HomotheticDistribution}
\end{proposition}

This is Açıkgöz's Proposition 7 at the level of measures. It is what lets the theorems be stated in
Aiyagari's own form, $K(r)=\E_r[a]$ at $w(r)$, while the household side is solved once at wage one
and the firm side reduces to $k/w=\alpha/((1-\alpha)(r+\delta))$, which needs no root
(\lean{aiyagari1994\_equilibriumRate\_unique\_wage}).

\subsection{Continuity of supply, for existence}\label{sec:cont}

Existence uses the intermediate value theorem, which needs capital supply continuous in the rate.
The rate is not a parameter of the household problem in the formal development but part of its
state, so that a family of problems indexed by the rate is one problem on a product space, and
continuity of the policy in the rate is joint continuity of one policy. Weak continuity of the
stationary distribution in the rate then follows from the Doeblin uniqueness (a limit of stationary
distributions is stationary, and there is only one), and aggregation is weakly continuous
because the asset coordinate is bounded and continuous (\lean{continuousOn\_capitalSupply}).

\subsection{What is never differentiated}\label{sec:deriv}

Every inequality above is derived without a derivative of the value or policy function. The Euler
inequalities come from perturbing the optimal plan by $\epsilon$ in one direction and letting
$\epsilon$ shrink: carrying the poorer household's plan forward, or the richer household's plan
back, gives one-sided bounds on the marginal value of assets, and the first-order condition is read
off the pair (\lean{euler\_le}, \lean{euler\_ge}; the construction is that of
\citealp{ClausenStrub2020}). Lipschitz constants are secant bounds; monotonicity in the rate is a
comparison of two plans. Where the development does establish differentiability of the value
function (\lean{IncomeFluctuationEnvelope}) it uses the Clausen--Strub sandwich and not the
Benveniste--Scheinkman argument, and the uniqueness proof does not depend on it.
